% Pursuing Stacks (cotes 134-2 to 134-8): archive cote 134-7, site page 110 (PDF page 111), typescript page 552.
% typescript. Draft compared: ps6.tex:5545-5617.
% AI-transcribed by Claude Opus 5.5 (claude-opus-5-5, adaptive thinking, effort high), Message Batches API, batch ps4 (msgbatch_012Aoc1NuskATPsyBeKSYK7J), 2026-11-12; not yet checked by eye.
\scanpage{134-7/110}
\mspage{552}

ves as ``the'' good object in terms of homological algebra and obstruction
theory in various geometric situations, \uline{only} in the case when the relati-%
ve scheme (X say) we are working with is \uline{smooth} over the base scheme S
- which in the present case amounts to saying (when $X=X_n$ is a component
of a semisimplicial ``bundle'') that X \xout{\unsure{its}} is indeed a flat vector bundle
over S. In more general cases, the work of André - Quillen - Illusie
tells us that the relevant object which replaces $\hand{\Omega}^1_{X/S}$ or it's dual $T_{X/S}$
is the relative tangent or cotangent \uline{complex} $L_\cdot^{X/S}$ or $L^\cdot_{X/S}$ , the second
being the dual of the other
\[ (15) \qquad L^\cdot_{X/S} = R\underline{Hom}_{\underline{O}_X}(L_\cdot^{X/S} , \underline{O}_X) , \]
these objects being viewed respectively as objects in the derived catego-%
ries $D_\cdot(\underline{O}_X)$ and $D^\cdot(\underline{O}_X)$ (deduced from chain and cochain complexes of
$\underline{O}_X$-Modules). When X is endowed with a section over S, the naïve diffe-%
rentials and codifferentials along this section should in the same way
be replaced by the co-Lie and Lie complexes
\[ (16) \qquad \hand{\ell}_\cdot(X/S,s) \quad\text{and}\quad \hand{\ell}^\cdot(X/S,s) \simeq R\underline{Hom}(\hand{\ell}_\cdot(X/S,s), \underline{O}_S) \quad , \]
where s is the given section, obtained from the previous complexes
by taking it's inverse images $Ls^*$ b\over{x}{y} s . As a matter of fact, the
chain complex $L_\cdot^{X/S}$ can be realized canonically, up to unique isomorphi-%
sm, via a semisimplicial Module
\[ L_*^{X/S} \]
on X, whose components are free $\underline{O}_X$-Modules. Accordingly, we get
(16) in terms of a well defined semisimplicial $\underline{O}_S$-Module\xout{\illegible{1}} \xout{\illegible{1}}
\[ (17) \qquad \hand{\ell}_*(X/S,s) \]
whose components are free - and as S=Spec(k), we may interprete this
more simply as a semisimplicial k-module with free components. \over{w}{W}hen we
apply this to the components $X_n$ of a semisimplicial bundle $X_*$ , we get
however the co-Lie invariants, in order to get the relevant Lie invaria-%
ants we'll have to take the duals
\[ (17') \qquad \hand{\ell}^*(X/S,s) = Hom_k(\hand{\ell}_*(X/S,s),k) \quad , \]
where X is any one among the $X_n$'s , and s it's marked section. Thus,
the ``corrected'' description of $L\hand{\Pi}_\cdot$ , by using the André-Quillen-Illusie
version of the ``Lie\hand{\unsure{$^*$}}-functor along a section'', would seem to be
\[ (18) \qquad L\hand{\Pi}_\cdot(X_*) = \hand{\ell}^*(\deleted{\illegible{1}}\, X_*/k , e_*) \ \hand{(?)}\hand{,} \]

\margin{110}
